\scanpage{193}
\begin{center}- \mspage{61} -\end{center}
\margin{IX}

dans ce cas
\[
(5.5.4)\qquad M=\operatorname{Hom}_{k\text{-gr}}(T'_{\unsure{o|0}},\mathbb G_{mk})\quad,
\]
et on aura alors
\[
\underline M=M_S\quad.
\]
Ceci posé, on aura évidemment, pour tout entier $n$, des isomorphismes
\[
(5.5.5)\qquad\underline M'_n\simeq D(T_n)\quad,\quad\underline M_n\simeq D(T'_n)\quad,
\]
où $T'_n$ est défini en termes de $A'$ comme $T_n$ en termes de $A$ (5.1.2), d'où on déduit des isomorphismes canoniques (où on pose $\mathbb Z_\ell(1)=T_\ell(\mathbb G_{mS})$) :
\[
(5.5.6)\qquad T_\ell(T_n)\simeq\check{\underline M}'_n\otimes_{\mathbb Z}\mathbb Z_\ell(1)\quad,\quad T_\ell(T'_n)\simeq\check{\underline M}_n\otimes_{\mathbb Z}\mathbb Z_\ell(1)\quad,
\]
et par passage à la limite sur $n$ des isomorphismes canoniques :
\[
(5.5.7)\qquad\begin{cases}
T_\ell(A^{\unsure{0|o}})^t=T_\ell(\hat T)\simeq\check{\underline M}'\otimes\mathbb Z_\ell(1)\\[1em]
T_\ell(A'^{\unsure{0|o}})^t=T_\ell(\hat T')\simeq\check{\underline M}\otimes\mathbb Z_\ell(1)\quad.
\end{cases}
\]
Dans le cas de la réduction semi-stable, on en conclut par dualité, en utilisant le théorème d'orthogonalité 5.2 et la relation (5.2.2), des isomorphismes canoniques :
\[
(5.5.8)\qquad\begin{cases}
T_\ell(A_K)/T_\ell(A_K)^f\simeq\underline M_K\otimes_{\mathbb Z}\mathbb Z_\ell\\[1em]
T_\ell(A')/T_\ell(A'_K)^f\simeq\underline M'_K\otimes_{\mathbb Z}\mathbb Z_\ell\quad.
\end{cases}
\]
On trouve de même, grâce à la même formule (5.2.2), des isomorphismes canoniques :
\[
(5.5.9)\qquad\begin{cases}
T_\ell(A_K)/T_\ell(A_K)^t\simeq D(T_\ell(A'_K)^f)\\[1em]
T_\ell(A'_K)/T_\ell(A'_K)^t\simeq D(T_\ell(A_K)^f)\quad,
\end{cases}
\]
\margin{379}
\margin{379}
\margin{192}
